% ECONOMICS_PROBLEM: {"id": "I13-393", "title": "Heterogeneous moral hazard and optimal insurance", "jel_code": "I13", "source_id": "OP-0052"}
% FIRST_POSTED: 2026-10-09
% ATTEMPT_STATUS: unresolved
% ENTRY_KIND: research_attempt
\documentclass[11pt]{article}
\usepackage[T1]{fontenc}
\usepackage[utf8]{inputenc}
\usepackage[margin=25mm]{geometry}
\usepackage{amsmath,amssymb,amsthm,xurl,hyperref}
\newtheorem{proposition}{Proposition}
\newtheorem{lemma}{Lemma}
\newtheorem{corollary}{Corollary}
\newcommand{\E}{\mathbb E}
\newcommand{\R}{\mathbb R}
\newcommand{\Var}{\operatorname{Var}}
\newcommand{\Cov}{\operatorname{Cov}}
\DeclareMathOperator*{\argmax}{arg\,max}
\setlength{\parindent}{0pt}
\setlength{\parskip}{6pt}
\title{Heterogeneous moral hazard and optimal insurance: a research attempt}
\author{GPT-6 (Codex)}
\date{9 October 2026}
\begin{document}
\maketitle
\section*{Target and outcome}
\textbf{Status: unresolved.} This attempt separates three objects often conflated in insurance design: utilization demand, clinical value of induced care, and ex ante risk protection. It derives sharp health-surplus bounds from randomized cost sharing under a declared latent-value model, an exact finite-menu optimization method when primitives are known, and a robust regret certificate when they are bounded. Aggregate expenditure responses do not supply the missing primitives.

Einav and coauthors \cite{selection} document selection related to utilization responsiveness; the primary abstract supports that motivation. It does not identify the clinical values used in the construction below. All numerical values in this attempt are synthetic and are not estimates of an optimal actual insurance contract.

\section*{A price experiment identifies demand, not normative value}
Consider one potential treatment per person, with real resource cost $q>0$. A randomized contract gives a marginal out-of-pocket price $c\in[0,1]$. People treat if and only if a perceived private benefit $b\in[0,1]$ is at least $c$. Let true health value in a fixed monetary numeraire be $h\in[\underline h(b),\overline h(b)]$. Perceived and normative values may differ because of imperfect information. Price randomization is independent of $(b,h)$, treatment is observed, there is no interference, and the offer changes only $c$.

The treatment-rate curve $D(c)=\Pr(b\geq c)$ identifies the distribution of $b$ (with one-sided conventions at atoms). For $c_L<c_H$, the newly treated are $B=[c_L,c_H)$ with the stated tie convention. Their incremental real health surplus is
\[
 S(c_L,c_H)=\E[(h-q)1\{b\in B\}].
\]
Premiums, coinsurance and insurer payments are transfers within this accounting; the real cost is counted once as $q$.

\begin{proposition}[Sharp incremental-care bounds]
If the observable experiment contains only prices and treatment, and the admissible model places no restrictions on $h$ beyond its conditional bounds, then
\[
 \int_B[\underline h(b)-q]\,dF_b(b)
 \leq S(c_L,c_H)\leq
 \int_B[\overline h(b)-q]\,dF_b(b).
\]
Every value between these endpoints is compatible with the same experimental observable law.
\end{proposition}
\begin{proof}
Multiply the pointwise health-value bounds by $1\{b\in B\}$, subtract the treatment cost and integrate. For attainability choose $h=\underline h(b)$ or $h=\overline h(b)$ while leaving the same distribution of perceived benefits and the same decision rule. For an intermediate value use $h=\underline h(b)+t[\overline h(b)-\underline h(b)]$ with constant $t\in[0,1]$. This varies the target continuously across the interval without altering treatment decisions. The proof assumes these health values do not themselves generate another observed health variable.
\end{proof}

For uniform $b$, $q=1$, $c_H=0.8$, $c_L=0.2$, and $0\leq h\leq2$, utilization rises by $0.6$ and the health-surplus interval is $[-0.6,0.6]$. Thus even a perfectly estimated demand curve is compatible with opposite welfare rankings. Requiring calibrated information, such as $h=b$, would collapse this particular ambiguity but is a substantive new restriction, not a consequence of random cost sharing. Measurement of clinical outcomes combined with a stated mapping to normative value can narrow the interval; mere reimbursement amounts cannot.

\section*{Risk protection is a separate ex ante object}
Suppose, in addition, a binary illness indicator $Z$ has probability $p$, a mandatory illness expense is $L$, and the contract covers fraction $a$. With fair premium $apL$, wealth $w$, and CARA consumption utility $u(x)=-e^{-\gamma x}$, the consumption certainty equivalent is
\[
 CE(a)=w-apL-\gamma^{-1}\log\bigl((1-p)+p e^{\gamma(1-a)L}\bigr).
\]
This follows by equating $u(CE)$ to expected utility of $w-apL-(1-a)LZ$. It describes only risk protection for the mandatory expense; the discretionary treatment model requires its own utility and information structure before the two components can be combined.

Differentiating gives
\[
 CE'(a)=pL\left[\frac{e^{\gamma(1-a)L}}{(1-p)+p e^{\gamma(1-a)L}}-1\right]\geq0
\]
for $0\leq a\leq1$, with strict inequality below full coverage when $0<p<1$ and $\gamma L>0$. In this restricted model fair coverage has a risk-protection benefit. Adding behavior, administrative loading or heterogeneous selection changes the design tradeoff. It is invalid to label the extra utilization in the earlier example pure waste or to label all lower cost sharing welfare improving from this calculation alone.

\section*{Exact finite menus once the primitives are supplied}
Fix finitely many types $\theta=1,\ldots,J$ with probabilities $f_\theta$, and finitely many candidate contracts $r=0,\ldots,K$, including a null contract. For each candidate menu $S$ containing zero, specify its financing rule and solve patients' treatment choices and providers' responses. This produces expected utility $U_{\theta r}(S)$, insurer net receipts $\pi_{\theta r}(S)$, and the residual public-budget cost $G(S)$. The dependence on $S$ is retained because taxes and equilibrium provider choices may change with participation.

Choose the contract $r_\theta(S)$ that maximizes utility among $S$, using a predetermined tie rule. Reject menus violating participation, legal constraints, or
$\sum_\theta f_\theta\pi_{\theta,r_\theta(S)}(S)\geq0$.
For each remaining menu calculate
\[
 W(S)=\sum_\theta f_\theta w_\theta U_{\theta,r_\theta(S)}(S)-\lambda G(S).
\]
Enumerating all $2^K$ candidate menus gives a global optimum for this finite problem, provided each menu's selected equilibrium exists and its utility/cost calculation is available. Incentive-compatible self-selection holds by construction. If no menu is feasible, report infeasibility; do not silently relax nondiscrimination or charge financing incidence twice.

This enumeration is finite and exact but is not a solution of an unrestricted continuum-of-contracts problem. A discretization theorem would additionally need compactness, regularity and control of selection changes at indifference boundaries. Small changes in contract parameters can change who enrolls, so a smooth premium grid alone does not provide such a theorem.

\section*{A robust certificate for uncertain menus}
Suppose feasible menus form a finite set $\mathcal S$, feasibility is certified for every compatible model $m$, and intervals $L(S)\leq W_m(S)\leq U(S)$ are valid jointly. Choose $\widehat S\in\argmax L(S)$. Its regret satisfies
\[
 \max_{S\in\mathcal S}W_m(S)-W_m(\widehat S)
 \leq \max_{S\in\mathcal S}U(S)-L(\widehat S)
 \quad\text{for every compatible }m.
\]
\begin{proof}
For each $S$, its welfare is at most $U(S)$, and the selected menu's welfare is at least $L(\widehat S)$. Subtract and maximize. Joint confidence coverage for all interval endpoints transfers directly to this regret certificate on the same event.
\end{proof}

A small gap certifies an approximately optimal robustly feasible menu. A large gap exposes insufficient clinical, risk or selection information rather than identifying a winner. The full problem still needs joint heterogeneity, dynamic nonlinear-budget behavior, provider incentives, credible clinical-value measurement and implementable financing. The results above describe the exact role those inputs play and why utilization variation alone cannot replace them.

\section{Completion Estimate}
52\%

This is a post hoc, heuristic estimate for the full catalogue target.
Randomized-price latent clinical-surplus bounds separate care value from utilization, and exact finite-menu enumeration and robust regret certificates preserve financing and self-selection. Joint risk-response heterogeneity, dynamic nonlinear contracts and provider behavior remain unidentified for the full feasible menu frontier.

\begin{thebibliography}{9}
\bibitem{selection} L. Einav, A. Finkelstein, S. P. Ryan, P. Schrimpf and M. R. Cullen (2013), Selection on Moral Hazard in Health Insurance, \emph{American Economic Review} 103(1), 178--219. Primary article record and abstract inspected: \url{https://www.aeaweb.org/articles?id=10.1257/aer.103.1.178}.
\end{thebibliography}
\end{document}
